\documentclass{article}
\usepackage[margin=1in]{geometry}
\usepackage{graphicx}
\usepackage{booktabs,multirow,amsmath,array}
\usepackage{float}
\begin{document}
\begin{table}[H]
\centering
\small
\setlength{\tabcolsep}{2.5pt}
\caption{Overview on MNIST (3 seeds). Network: 784 inputs, one hidden layer of 128 leaky integrate-and-fire (LIF) spiking neurons, and 10 outputs; each image is presented for 25 time steps. Comparison 1: backprop vs ours (dense). Comparison 2: backprop with our energy term vs ours (sparse). Bold: better. $^\dagger$Within one std.\\}
\label{tab:overview}
\resizebox{\textwidth}{!}{%
\begin{tabular}{>{\raggedright\arraybackslash}p{2.1cm}>{\raggedright\arraybackslash}p{6cm}cccc}
\toprule
Category & Metric & Backprop & Ours (dense) & Backprop + energy & Ours (sparse) \\
\midrule
\multirow{2}{*}{\parbox[c]{2.1cm}{\raggedright\textit{Accuracy}}} & Test accuracy (\%)$^\dagger$ & $\boldsymbol{97.63 \pm 0.05}$ & $97.59 \pm 0.09$ & $\boldsymbol{97.12 \pm 0.07}$ & $97.06 \pm 0.15$ \\
 & ECE (calibration) & $\boldsymbol{0.0034 \pm 0.0002}$ & $0.0038 \pm 0.0001$ & $\boldsymbol{0.0044 \pm 0.0010}$ & $0.0087 \pm 0.0007$ \\
\midrule
\multirow{5}{*}{\parbox[c]{2.1cm}{\raggedright\textit{Energy}}} & Synaptic operations per image, network$^\dagger$ & $11339 \pm 243$ & $\boldsymbol{3858 \pm 474}$ & $525 \pm 120$ & $\boldsymbol{472 \pm 15}$ \\
 & Synaptic operations per image, total$^\dagger$ & $343802 \pm 255$ & $\boldsymbol{336322 \pm 466}$ & $332989 \pm 128$ & $\boldsymbol{332936 \pm 24}$ \\
 & GPU energy per training step (mJ) & $991 \pm 19$ & $\boldsymbol{768 \pm 11}$ & $1078 \pm 4$ & $\boldsymbol{682 \pm 3}$ \\
 & Total training GPU energy (kJ) & $17.78 \pm 1.23$ & $\boldsymbol{14.10 \pm 0.15}$ & $23.10 \pm 7.57$ & $\boldsymbol{14.21 \pm 1.43}$ \\
 & Inference GPU energy per image (mJ)$^\dagger$ & $\boldsymbol{1.48 \pm 0.13}$ & $1.49 \pm 0.03$ & $\boldsymbol{1.42 \pm 0.03}$ & $1.44 \pm 0.06$ \\
\midrule
\multirow{5}{*}{\parbox[c]{2.1cm}{\raggedright\textit{Efficiency}}} & Training step (ms) & $7.64 \pm 0.05$ & $\boldsymbol{5.53 \pm 0.16}$ & $9.21 \pm 0.01$ & $\boldsymbol{5.37 \pm 0.03}$ \\
 & Training time per epoch (s) & $4.57 \pm 0.19$ & $\boldsymbol{3.82 \pm 0.24}$ & $5.74 \pm 1.43$ & $\boldsymbol{3.73 \pm 0.13}$ \\
 & Epochs to 95\% val. accuracy$^\dagger$ & $2.0 \pm 1.0$ & $\boldsymbol{1.7 \pm 0.6}$ & $3.0 \pm 0.0$ & $\boldsymbol{1.3 \pm 0.6}$ \\
 & Training time to 95\% val. accuracy (s)$^\dagger$ & $9.4 \pm 4.3$ & $\boldsymbol{6.7 \pm 1.9}$ & $17.8 \pm 5.0$ & $\boldsymbol{5.8 \pm 1.8}$ \\
 & Peak GPU memory, training step (MiB) & $\boldsymbol{230.9 \pm 0.0}$ & $271.6 \pm 0.0$ & $\boldsymbol{234.1 \pm 0.0}$ & $248.5 \pm 0.0$ \\
\midrule
\multirow{3}{*}{\parbox[c]{2.1cm}{\raggedright\textit{Explainability and robustness}}} & Faithfulness margin, 25\% removed & $0.04 \pm 0.02$ & $\boldsymbol{0.40 \pm 0.15}$ & $\boldsymbol{0.41 \pm 0.07}$ & $0.10 \pm 0.02$ \\
 & Credit vs. neuron removal (Spearman) & $0.57 \pm 0.03$ & $\boldsymbol{0.81 \pm 0.11}$ & $\boldsymbol{0.93 \pm 0.03}$ & $0.84 \pm 0.04$ \\
 & Accuracy after random 25\% removal$^\dagger$ & $\boldsymbol{92.8 \pm 1.0}$ & $81.3 \pm 6.1$ & $89.5 \pm 2.2$ & $\boldsymbol{91.4 \pm 0.4}$ \\
\midrule
\textit{Backprop. Check} & Autograd calls in training (per seed) & 10{,}700 & 0 & 10{,}700 & 0 \\
\bottomrule
\end{tabular}%
}
\end{table}

{
\noindent\textbf{Overview of the four fronts}\par\smallskip
\noindent\textbf{Accuracy:}
\begin{itemize}
  \item \textbf{Test accuracy (\%):} The share of the 10,000 test images classified correctly, using the checkpoint with the best validation accuracy. The prediction is the read-out neuron with the most accumulated membrane potential over the 25 time steps. Higher is better.
  \item \textbf{ECE (calibration):} Expected calibration error. Test predictions are grouped into 15 bins by confidence; in each bin, the gap between average confidence and actual accuracy is taken, and the gaps are averaged, weighted by bin size. It measures whether a model's confidence matches how often it is right (0 = perfect). Probabilities are temperature-scaled on the validation set first, so both models are compared at their best calibration. Lower is better.
\end{itemize}

\noindent\textbf{Energy:}
\begin{itemize}
  \item \textbf{Synaptic operations per image, network:} How many synapses are activated by hidden-layer spikes while classifying one test image: a spike of a neuron with $n$ outgoing synapses counts $n$. It is the work a neuromorphic chip would do, which only computes when a neuron fires. Lower is better.
  \item \textbf{Synaptic operations per image, total:} The same, including the input layer's spikes. The input is identical for both models and dominates this count in a network with one hidden layer, so the difference here is small even when the network count differs severalfold. Lower is better.
  \item \textbf{GPU energy per training step (mJ):} Energy the GPU draws for one training step (forward pass plus the model's own update; for backprop this includes its backward pass), read from the GPU's energy counter in the controlled benchmark. Lower is better.
  \item \textbf{Total training GPU energy (kJ):} GPU energy over the whole 50-epoch training run, including validation passes. It reflects the total compute spent to train a model. Lower is better.
  \item \textbf{Inference GPU energy per image (mJ):} GPU energy to classify one test image, from the controlled benchmark. On a GPU, which computes densely, sparsity does not lower this; the synaptic-operation counts show the hardware-independent difference. Lower is better.
\end{itemize}

\noindent\textbf{Efficiency:}
\begin{itemize}
  \item \textbf{Training step (ms):} Time for one training step on a batch of 256 images, from the controlled benchmark (same preloaded data, models interleaved, repeated). Lower is better.
  \item \textbf{Training time per epoch (s):} Time for one pass over the 55,000 training images during the actual training runs, excluding validation. Lower is better.
  \item \textbf{Epochs to 95\% val.\ accuracy:} After every epoch, the model is evaluated on the 5,000 validation images. This is the first epoch at which that accuracy reaches 95\%, averaged over the 3 seeds. It measures how many passes over the training data a model needs to become useful. Lower is better.
  \item \textbf{Training time to 95\% val.\ accuracy (s):} The training time, in seconds, from the start until the end of that epoch: the per-epoch training times added up. For example, a model that reaches 95\% at the end of its second epoch, with epochs of 3.8\,s, scores 7.6\,s. The validation checks between epochs are not added, because they only measure the model and do not train it. It combines how many epochs are needed with how long each one takes. Lower is better.
  \item \textbf{Peak GPU memory, training step (MiB):} The most GPU memory in use during a training step in the controlled benchmark, including the preloaded batches. Our method keeps information about every time step of the window (for its local updates), which costs more memory than backprop's stored activations. Lower is better.
\end{itemize}

\noindent\textbf{Explainability and robustness:}
\begin{itemize}
  \item \textbf{Faithfulness margin, 25\% removed:} Each model assigns a credit score to each of its 128 hidden neurons. The margin is test accuracy after removing 32 random neurons minus test accuracy after removing the 32 top-credited ones (best of the three credit scores). Positive means the credit identifies the neurons that matter. It also rewards networks that concentrate their computation in a few neurons. Higher is better.
  \item \textbf{Credit vs.\ single-neuron removal (Spearman):} Each hidden neuron is removed on its own and the resulting increase in test loss is measured. The score is the rank correlation between those measured effects and the credit scores (1 = the credit orders the neurons exactly by their real importance). Unlike the margin, it does not depend on how concentrated the network is. Higher is better.
  \item \textbf{Accuracy after removing a random 25\% (\%):} Test accuracy after silencing 32 randomly chosen hidden neurons, averaged over five random draws. It measures robustness to losing neurons (for example, hardware faults). Higher is better.
\end{itemize}

\noindent\textbf{Backprop. check:}
\begin{itemize}
  \item \textbf{Autograd calls in training (per seed):} How many times PyTorch's backpropagation engine was called during training. Backprop calls it once per batch (214 batches $\times$ 50 epochs = 10,700). Ours must be 0; this verifies that our method never uses backprop.
\end{itemize}
}

\end{document}
